
\documentclass{../concepts/prepareprotokol} %všechny balíčky pro protokol + titulní strana
\usepackage{graphicx}
\usepackage{booktabs}
\usepackage{amsmath}
\usepackage{newfloat}
\usepackage{caption}
\usepackage{siunitx}
\sisetup{inter-unit-product =$\cdot$}
\sisetup{separate-uncertainty=true}

\praktikum{II} %číslo praktika (I, II nebo III)
\autor{David Němec} %jméno autora
\datum{10.\,11.\,2025} %datum měření
\cislo{4} %číslo úlohy
\nazev{Měření malých odporů} %název úlohy

\DeclareFloatingEnvironment[fileext=los, placement={htbp}, name=Obrázek]{schema}

\newcommand{\nobars}{
Chybové úsečky nebyly pro svou malou velikost pro přehlednost vykresleny.}


\begin{document}
\renewcommand{\figurename}{Graf}
\renewcommand{\refname}{Literatura}

\maketitle

% ========================== PRACOVNÍ ÚKOL =========================
\section{Pracovní úkol}
\begin{enumerate}
    \item Změřte průměry a délky drátů na pracovní desce.

    \item Změřte odpory ve čtyřbodovém zapojení pomocí multimetru KEITHLEY 2010.
    
    \item Změřte odpor drátů Wheatstoneovým a Thomsonovým můstkem Metra - MTW. Vysvětlete rozdíly ve výsledcích měření. Současně určete odpor přívodních vodičů a odpor na svorkách v případě měření Wheatstoneovým můstkem.
    
    \item Určete měrný odpor jednotlivých vzorků i s příslušnou chybou výsledku. Stanovené hodnoty porovnejte s hodnotami uváděnými v tabulkách.

\end{enumerate}

% ============================= TEORIE ===============================
\section{Teorie}

\subsection{Čtyřbodové zapojení rezistoru}

Pro přesnější měření odporu rezistoru lze využít čtyřbodového zapojení podle schématu na obrázku \ref{sch:four}. Pokud je voltmetr připojen k napěťovým svorkám $b$ a $b'$, nebude měření ovlivněno spádem napětí na přívodních vodičích mezi rezistorem a proudovými svorkami $a$ a $a'$. Za předpokladu nízkých hodnot $R$ oproti odporu voltmetru poteče přívodními vodiči k napěťovým svorkám zanedbatelný proudu a ani odpor těchto vodičů se neuplatní \cite{studijni}.

\subsection{Wheastonův můstek}

Nízké hodnoty odporu studovaného rezistoru lze měřit pomocí Wheastonova můstku. Jeho schéma zapojení je na obrázku \ref{sch:wheast}. Je-li můstek v rovnováze, tj. galvanometrem G neteče proud, platí pro poměry jednotlivých odporů $R_1$, $R_2$, $R_3$ a $R_4$ vztah \cite{studijni}
\begin{equation}
    \frac{R_1}{R_2} = \frac{R_3}{R_4} \, .
\end{equation}
Zafixováním poměru $R_3$\,/\,$R_4$ na vysoké hodnotě a postupným měněním odporu $R_2$ tak, aby výchylka galvanometru byla nulová, lze pak snadno změřit zkoumaný nízký odpor $R_1$. V jeho hodnotě však bude v důsledku tohoto zapojení započten i odpor přívodních drátů.

\subsection{Thomsonův můstek}

Přesnější měření lze provést pomocí Thomsonova můstku, jehož schéma zapojení je na obrázku \ref{sch:four}. Pokud bude opět můstek v rovnováze (galvanometrem G nepoteče proud) a pro odpory $P$, $Q$, $p$ a $q$ bude platit
\begin{equation}
    \frac{p}{P} = \frac{q}{Q} \, ,
\end{equation}
lze zkoumaný odpor $R_X$ určit jednoduše jako \cite{studijni}
\begin{equation}
    R_X = \frac{Q}{P} R_N \, .
\end{equation}
Při čtyřbodovém zapojení odporů $R_N$ a $R_X$ již nebude měření ovlivněno odporem jejich přívodních drátů. Za předpokladu nízkých hodnot $R_N$ a $R_X$ oproti odporům $P$, $Q$, $p$ a $q$, budou také nízké proudy $I$ a $i$ oproti proudu $I_0$ a odpory přívodních vodičů k můstku se díky tomu také v měření neuplatní \cite{studijni}.

\subsection{Měrný odpor}

Měrný odpor $\rho$ drátu průřezu $S$ a délky $l$ je definován vztahem \cite{studijni}
\begin{equation}
    \rho = \frac{R S}{l} \, ,
\end{equation}
kde $R$ je celkový odpor drátu. Je-li příčný řez drátu kruhový, lze průřez $S$ vypočítat jako
\begin{equation}
    S = \frac{\pi D^2}{4} \, ,
\end{equation}
kde $D$ je průměr drátu.


\subsection{Metoda měření}

Nejprve byly zjištěny rozměry zkoumaných drátů - průměr a délka, přičemž délka byla měřena mezi napěťovými kontakty na obou koncích drátu. Pro větší přesnost byl průměr měřen na několika místech drátu.

Při měření multimetrem KEITHLEY 2010 bylo použito čtyřbodové zapojení rezistoru (obr. \ref{sch:four}), napěťové a proudové kontakty byly propojeny se vstupními svorkami multimetru podle pokynů v \cite{pokyny}.

Při měření odporu pomocí Wheastonova a Thomsonova můstku byly nastaveny hodnoty zbylých odporů ($R_3$ a $R_4$ pro Wheastonův můstek a $P$ pro Thomsonův můstek) tak, aby hodnota na odporové dekádě v jednotkách \si{\ohm} odpovídala hodnotě zkoumaného odporu v jednotkách \si{\milli \ohm} (viz \ref{wheaston} a \ref{thomson}). Jako první odhad optimální hodnoty odporu na odporové dekádě posloužila hodnota odporu určená pomocí multimetru KEITHLEY 2010. Otáčením postupně jemnějších koleček dekády byla vyladěna výsledná hodnota tak, aby byl můstek v rovnováze, tedy aby galvanometrem netekl proud a jeho výchylka byla nulová.

\begin{schema}[htbp]
    \centering
    \includegraphics[width=0.5\linewidth]{uloha4/fourpoints.png}
    \caption{Čtyřbodové zapojení rezistoru \cite{studijni}}
    \label{sch:four}
\end{schema}

\begin{schema}[htbp]
    \centering
    \includegraphics[width=0.5\linewidth]{uloha4/wheaston.png}
    \caption{Wheastonův můstek \cite{studijni}}
    \label{sch:wheast}
\end{schema}

\begin{schema}[htbp]
    \centering
    \includegraphics[width=0.5\linewidth]{uloha4/thomson.png}
    \caption{Thomsonův můstek \cite{studijni}}
    \label{sch:thomson}
\end{schema}


\subsection{Měřicí přístroje a jejich nejistoty}
\begin{enumerate}
    \item Mikrometr s nejmenším dílkem 0.01 mm:
    
    Mikrometrem byly měřeny průměry jednotlivých drátů z různých materiálů. Nejistotu naměřené hodnoty odhaduji jako velikost nejmenšího dílku, tedy 0.01 mm.

    \item Pásový metr s nejmenším dílkem 1 mm:

    Metrem byly měřeny délky jednotlivých drátů. Nejistotu naměřené hodnoty odhaduji jako velikost nejmenšího dílku, tedy 1 mm. Některé dráty však nebyly mezi svorkami napnuty úplně dokonale, proto je nejistota jejich délky větší (viz \ref{rozmery}).

    \item Multimetr KEITHLEY 2010:

    Multimetrem se měřily odpory drátů ve čtyřbodovém zapojení. Z manuálu k použití přístroje v příloze v \cite{pokyny} je nejistota měření odporu (pro rozsah \SI{10}{\ohm}, na kterém bylo měřeno) 60 ppm z naměřené hodnoty + 9 ppm z rozsahu.

    \item Odporová dekáda:

    Nastavováním odporu na odporové dekádě byla hledána rovnováha na Wheastonově a Thomsonově můstku. Nejistota nastaveného odporu byla odhadnuta jako největší řád, který již neměl vliv na výchylku galvanometru (tj. výchylka byla nulová a nijak se něměnila ani při otáčení příslušného kolečka odporové dekády). Nejistota dekády uvedená na její přední straně je zanedbatelná vůči takto odhadnuté nejistotě pro všechny řády kromě řádu \SI{1000}{\ohm}. Jeho nejistota je 0.02\,\% z nastavené hodnoty.

    \item Termohygrobarometr Commeter C4130:

    termohygrobarometrem byla zjištěna teplota okolí během měření, aby experimentálně určené hodnoty měrného odporu mohly být srovnány s tabulkovými. Nejistota teploty byla určena z manuálu přístroje jako \SI{0.4}{\celsius}.

\end{enumerate}

% ======================== VÝSLEDKY MĚŘENÍ ===========================
\section{Výsledky měření}

Průměry jednotlivých drátů byly změřeny mikrometrem na deseti místech a z nich byl vypočtena střední hodnota (aritmetický průměr). Naměřené hodnoty a jejich střední hodnoty shrnuje tabulka \ref{tab:prumer}. Nejistota středních hodnot byla vypočtena jako 
\begin{equation}
    \sigma_{\bar{D}} = \sqrt{\frac{1}{9} \sum_{i=1}^{10} (\bar{D} - D_i)^2 + \sigma_{mic}^2} \,
\end{equation}
kde $\bar{D}$ je střední hodnota průměru drátu, $D_i$ je i-tá hodnota průměru a $\sigma_{mic}=\SI{0.01}{\milli\metre}$ je nejistota měřidla.

Vypočtené střední hodnoty průměrů drátů a jejich délky ukazuje tabulka \ref{tab:delky}. Nejistota určení délky byla u většiny drátů vzata jako nejistota pásového metru. Měděný drát však zjevně nebyl mezi svorkami dokonale napnutý, proto jeho délka i její nejistota byly odhadnuty jako o něco větší, než byla přímá vzdálenost mezi svorkami.

\begin{table}[!hbtp]
\centering
\caption{Změřené průměry drátů}
\begin{tabular}{cccccc}
 \toprule
 wolfram & měď & kanthal & železo & mosaz & chromnikl \\
 \midrule
 
& & & & & \\
 & & & & & \\
 & & & & & \\
 & & & & & \\
 & & & & & \\
 & & & & & \\
 & & & & & \\
 & & & & & \\
 & & & & & \\
 & & & & & \\
 
 \midrule

 & & & & & \\
 
 \bottomrule
\end{tabular}
\label{tab:prumer}
\end{table}


% ============================ DISKUZE ===============================
\section{Diskuze}


% ============================= ZÁVĚR ================================
\section{Závěr}


% =========================== LITERATURA =============================
\begin{thebibliography}{9}

\bibitem{studijni} Kolektiv ZFP KVOF MFF UK: Měření malých odporů [online]. [cit. 17.11.2025], dostupné z~https://physics.mff.cuni.cz/vyuka/zfp/\_media/zadani/texty/txt\_204.pdf

\bibitem{pokyny} Kolektiv ZFP KVOF MFF UK: Měření malých odporů - pokyny k měření [online]. [cit. 17.11.2025], dostupné
z~https://physics.mff.cuni.cz/vyuka/zfp/\_media/zadani/pokyny/mereni\_204.pdf

\bibitem{englich} J. Englich: Úvod do praktické fyziky I: Zpracování výsledků měření. 1. vyd. Praha: Matfyzpress, 2006

\bibitem{mikulcak} Mikulčák a kolektiv: Matematické, fyzikální a chemické tabulky pro střední školy: relativní permitivita, Praha: Státní pedagogické nakladatelství, n. p., 1988

\end{thebibliography}

\end{document}